\section{A self-contained analytic foundation for the zero results}
\label{sec:zero-preliminaries}

The new zero-count arguments require a bound valid on the entire positive
axis, including multiple zeros. We give the necessary proof here. This
section reproduces the mechanism already established in the baseline;
the subsequent classical and deformation results are the continuation.

For $n\ge0$, $k\ge1$, and $0\le\rho\le1$, set
\begin{equation}\label{eq:unified-D}
 D_{n,k}^{\rho}(a)=\sum_{m\ge0}\rho^m
       \frac{d^k}{da^k}\frac{\log^n(a+m)}{a+m}.
\end{equation}
The summands are $O(m^{-k-1}(1+\log^n m))$ on compact positive
$a$ intervals. Thus this series and all its fixed $a$ derivatives converge
normally, uniformly in $\rho\in[0,1]$. For $\rho<1$ it equals the
$k$th derivative of
$(-1)^n\partial_s^n\Phi(\rho,s,a)|_{s=1}$.

At $\rho=1$, differentiating the Hurwitz identity
\[
 \partial_a^k\zeta(s,a)=(-1)^k(s)_k\zeta(s+k,a)
\]
in a right half-plane and continuing to $s=1$ proves
$D_{n,k}^1=\gamma_n^{(k)}$. More explicitly, if
\[
 P_k(z)=\prod_{j=1}^k(1+z/j)=\sum_i e_i(k)z^i,
\]
then coefficient comparison at $s=1+z$ gives
\begin{equation}\label{eq:master-coefficients}
 D_{n,k}^{\rho}(a)=(-1)^{n+k}k!n!
 \sum_{i+j=n}\frac{e_i(k)}{j!}
     \partial_s^j\Phi(\rho,k+1,a).
\end{equation}
At $\rho=1$, the last $\Phi$ denotes $\zeta$. The pole in
\eqref{eq:laurent-convention} is independent of $a$, which is why no pole
term remains after a positive number of parameter derivatives. This is
the harmonic-coordinate version of the all-order formula treated by
Coffey \cite{Coffey}.

\subsection{The positive weight and the reciprocal-gamma polynomial}

Define monic Appell polynomials by
\begin{equation}\label{eq:Q-generating}
 \frac{e^{xz}}{\Gamma(1+z)}=\sum_{n\ge0}Q_n(x)\frac{z^n}{n!}.
\end{equation}
For example, with $y=x+\gamma$,
\[
 Q_1=y,\qquad Q_2=y^2-\zeta(2),\qquad
 Q_3=y^3-3\zeta(2)y+2\zeta(3).
\]
The Mellin formula and $(1+z)_k/\Gamma(k+1+z)=1/\Gamma(1+z)$
give
\begin{equation}\label{eq:positive-Laplace}
 (-1)^{n+k}D_{n,k}^{\rho}(a)
 =\int_0^\infty\frac{t^k e^{-at}}{1-\rho e^{-t}}Q_n(\log t)\dd t.
\end{equation}
To justify coefficient extraction and differentiation, use
\[
 1\le\frac1{1-\rho e^{-t}}\le\frac1{1-e^{-t}}\le1+\frac1t.
\]
Near zero the integrable majorants are
$t^{k-1-\epsilon}(1+|\log t|^n)$ with $0<\epsilon<k$;
at infinity a positive lower bound on $a$ supplies exponential decay.
Any fixed number of $a$ derivatives only multiplies by a power of $t$.
These estimates also prove continuity at $\rho=1$.

\begin{lemma}[Root preservation and strictness]\label{lem:root-operator}
If a real polynomial $p$ has only real roots and $c\ne0$ is real, then
$p+cp'$ has only real roots. A root of multiplicity $m\ge2$ is inherited
with multiplicity $m-1$; all other roots are simple. Consequently any
sequence of $n-1$ such operators makes a degree-$n$ real-rooted polynomial
strictly real-rooted.
\end{lemma}
\begin{proof}
Away from the distinct roots $r_i$ of $p$, of multiplicities $m_i$, the
new roots satisfy
\[
 \frac{p'}p=\sum_i\frac{m_i}{x-r_i}=-\frac1c.
\]
The left side has strictly negative derivative on each complementary
interval. Its endpoint limits give one simple solution in each gap and
one in the appropriate exterior interval. Factoring at $r_i$ gives the
inherited multiplicity, and the degree accounts for all roots.
Repeated application reduces every possible multiplicity without
introducing new multiple roots.
\end{proof}

\begin{proposition}\label{prop:Q-simple}
For each $n\ge1$, $Q_n$ has $n$ distinct real roots. The zeros of
$Q_{n-1}$ strictly interlace them.
\end{proposition}
\begin{proof}
Use the locally uniform Weierstrass product \cite{DLMFGamma}
\[
 \frac1{\Gamma(1+z)}=e^{\gamma z}
       \prod_{m\ge1}(1+z/m)e^{-z/m}.
\]
For a finite product, extracting $n![z^n]$ after multiplying by $e^{xz}$
amounts to applying real translations and the operators
$1+\partial_x/m$ to $x^n$. Lemma~\ref{lem:root-operator} proves real
rootedness of every such polynomial. Coefficient convergence preserves
real rootedness for monic polynomials of fixed degree.

Simplicity requires an additional step: a limit of polynomials with simple
roots need not have simple roots. Remove the first $n-1$ factors from the
infinite product. The same argument makes the polynomial associated with
the tail real-rooted. Restore those $n-1$ factors; their translations
commute with the differential operators, and the lemma makes every root
simple. Finally $Q_n'=nQ_{n-1}$ gives strict interlacing by Rolle's theorem.
\end{proof}

\subsection{The global bound counts multiplicities}

\begin{lemma}[Laplace variation bound]\label{lem:laplace-variation}
Suppose $g(t)$ has exactly $N$ sign changes on $(0,\infty)$, with a
fixed nonzero sign between successive sign-change points up to null sets.
Suppose its exponentially weighted polynomial moments are absolutely
integrable for every positive Laplace parameter. Then
$G(a)=\int_0^\infty e^{-at}g(t)\dd t$ is nonzero as a function and has
at most $N$ zeros on $(0,\infty)$, counted with multiplicity.
\end{lemma}
\begin{proof}
For $N=0$ the integral has fixed strict sign. For $N>0$, choose a
sign-change point $t_0$. The function $(t_0-t)g(t)$ has one fewer sign
change, and its transform is
\[
 (\partial_a+t_0)G(a)
 =e^{-t_0a}\frac{d}{da}\bigl(e^{t_0a}G(a)\bigr).
\]
Induction bounds its zeros by $N-1$. If any finite collection of zeros
of $G$ has total multiplicity $M$, Rolle's theorem with multiplicity
gives at least $M-1$ zeros of the displayed derivative, including the
inherited endpoint multiplicities. Hence $M\le N$. This applies to every
finite collection. The induction also proves nonvanishing: an identically
zero $G$ would have an identically zero transformed derivative.
\end{proof}

\begin{theorem}[Baseline global zero theorem, reproduced]
\label{thm:global-zero-bound}
For every $n\ge0$, $k\ge1$, and $0\le\rho\le1$,
$D_{n,k}^{\rho}$ has at most $n$ positive zeros, counted with multiplicity.
\end{theorem}
\begin{proof}
In \eqref{eq:positive-Laplace}, the factor $t^k/(1-\rho e^{-t})$ is
strictly positive. Proposition~\ref{prop:Q-simple} says that
$Q_n(\log t)$ has exactly $n$ simple sign changes. The domination already
proved supplies the moment hypotheses of Lemma~\ref{lem:laplace-variation}.
That lemma proves the conclusion, including nonvanishing.
\end{proof}

This bound will serve two distinct purposes. At a derivative order with
$n$ certified sign changes it proves completeness and simplicity. At a
parameter where a double zero is known to exist, it can rule out higher
degeneracy once the other zeros use the remaining available multiplicity.
The latter use is what makes the Lerch bifurcation proof rigorous.
